% Folder 93, batch 01 — archivists' pages 1–20.
% Pass: Opus 5.5 (claude-opus-5-5), 2026-10-10 — first pass, unchecked against the pages by a human.
%
% Pages 2–19 are a typescript by P. Deligne (« Courbes elliptiques : formulaire,
% d'après J. Tate »), third-party work under RIGHTS.md: not transcribed, one note
% per page. The handwriting on the copy (black, green and blue ink, pp. 2–3,
% 4, 5, 9, 10, 14, 15, 17, 18, 20) is not taken to be Grothendieck's: the same
% hand writes on p. 2 « G? = correction de Grothendieck », i.e. it attributes,
% with a query, certain corrections to him in the third person. It is therefore
% recorded, not transcribed. Page 1 is blank and omitted.
\documentclass[11pt,a4paper]{article}
\input{../preamble/grothendieck}

\folder{93}
\batch{1}
\pages{1}{20}
\foldertitle{Connexions de Gauss-Manin et équations de Picard-Fuchs. Opération de Cartier : copies de tapuscrit annoté (s.d.), tiré à part (1981), notes manuscrites (s.d.).}
\dating{[1972]-1981}
\watermark{Édition de démonstration}

\begin{document}

\section{Courbes elliptiques : formulaire (tapuscrit d'un tiers)}

\note{Les pages 2 à 19 sont une copie d'un tapuscrit qui n'est pas de Grothendieck : « Courbes elliptiques : formulaire, d'après J. Tate, mis au goût du jour par P. Deligne », paginé 1 à 16, sans date. Texte d'un tiers : il n'est pas transcrit ici. La copie porte des annotations manuscrites (encres noire, verte et bleue) d'une main qui écrit en tête de la p. 2 « G? = correction de Grothendieck » : elle parle de lui à la troisième personne et n'est donc pas prise pour la sienne ; ces annotations sont signalées, non transcrites. Les quelques corrections que cette main marque « G? » sont indiquées là où elles se trouvent ; leur attribution à Grothendieck est la sienne, et elle est elle-même interrogative.}

\page{2}
\note{Tapuscrit de P. Deligne d'après J. Tate, s.d., p. 1 : titre, § 1 « Formules générales » (courbes de genre 1 sur une base $S$, types de fibres elliptique, multiplicatif, additif ; faisceau $\omega = e^{*}\Omega^{1}_{E/S}$). En tête, de la main annotatrice : la légende « G? = correction de Grothendieck », et un renvoi à une version plus complète (Lecture Notes 476, Anvers IV, et l'exposé de Tate dans le même volume). En marge, « G? » en regard de « arithmétique », souligné à l'encre bleue.}

\page{3}
\note{Feuillet manuscrit de la main annotatrice, à l'encre verte, intitulé « Formules utiles » : relation $1728\Delta = c_4^3 - c_6^2$, invariant $j$, formes $c_4$, $c_6$, $\Delta$ comme sections de puissances de $\omega$, développements de $c_4$, $c_6$ pour la courbe de Tate, invariant de Hasse, courbe de Legendre. Écriture d'un tiers : non transcrite.}

\page{4}
\note{Tapuscrit de Deligne, p. 2 : filtration de $p_{*}\mathcal{O}(ne)$, formule (1.1), faisceau des différentielles régulières $\Omega^{\mathrm{reg}}_{E/S}$. Deux renvois marginaux « réf » de la main annotatrice, à l'encre bleue.}

\page{5}
\note{Tapuscrit de Deligne, p. 3 : différentielles invariantes, changements de coordonnées (1.2), équation de Weierstrass (1.3). En marge, « G? » en regard de la ligne où le signe $\otimes$ de l'exposant $\otimes{-2}$ est repassé à l'encre bleue.}

\page{6}
\note{Tapuscrit de Deligne, p. 4 : forme différentielle invariante, quantités $b_2, b_4, b_6, b_8, d_6, d_8$ (1.4). Sans annotation manuscrite.}

\page{7}
\note{Tapuscrit de Deligne, p. 5 : formules (1.5) pour $c_4$, $c_6$, $\Delta$, et formules de changement de variables (1.6), (1.7). Sans annotation manuscrite.}

\page{8}
\note{Tapuscrit de Deligne, p. 6 : (1.8) ; § 2 « Cas où 2 et 3 sont inversibles », forme de Weierstrass, proposition 2.5 (schéma de modules au-dessus de $\operatorname{Spec}\mathbb{Z}[2^{-1},3^{-1}]$). Sans annotation manuscrite.}

\page{9}
\note{Tapuscrit de Deligne, p. 7 : fin de la proposition 2.5, relation (2.6) ; § 3 « Cas où 2 et $c_4$ sont inversibles ». En marge, de la main annotatrice à l'encre bleue, un point d'interrogation et « section $e$ ».}

\page{10}
\note{Tapuscrit de Deligne, p. 8 : formules (3.3), (3.4), proposition 3.5. Même annotation marginale « ? section $e$ » à l'encre bleue.}

\page{11}
\note{Tapuscrit de Deligne, p. 9 : § 4 « Cas où 2 est nilpotent et $c_4$ inversible » ; début du § 5 « Cas où $c_4 = 0$, $p = 2$ ou $p = 3$ ». Sans annotation manuscrite.}

\page{12}
\note{Tapuscrit de Deligne, p. 10 : suite du § 5, caractéristiques 3 et 2. Sans annotation manuscrite.}

\page{13}
\note{Tapuscrit de Deligne, p. 11 : § 6 « Applications », proposition 6.1 (types de fibres selon $\Delta$ et $c_4$). Sans annotation manuscrite.}

\page{14}
\note{Tapuscrit de Deligne, p. 12 : invariant $j$ (6.2), proposition 6.3 sur $\mathrm{Isom}(E,E')$. L'indice $S$ est ajouté à l'encre bleue sous les trois $\mathrm{Isom}$ de l'énoncé.}

\page{15}
\note{Tapuscrit de Deligne, p. 13 : proposition 6.4 et sa preuve. Une parenthèse complétée à l'encre bleue dans « $1/(j-1728)$ ».}

\page{16}
\note{Tapuscrit de Deligne, p. 14 : suite de la preuve ; formes types (6.5), (6.6). Sans annotation manuscrite.}

\page{17}
\note{Verso resté blanc du tapuscrit, portant seulement des notes de la main annotatrice à l'encre bleue (« Question connexe » sur le schéma formel des modules d'une courbe elliptique sur $W(k)$ et le quotient par $\mathrm{Aut}(E)$, cas $j = 0$ et $j = 12^3$). Écriture d'un tiers : non transcrite.}

\page{18}
\note{Tapuscrit de Deligne, p. 15 : formes types (6.7), (6.8) ; proposition 6.9 sur les automorphismes. En marge de (I), à l'encre bleue, une réserve de la main annotatrice ; un point d'interrogation en regard de (V).}

\page{19}
\note{Tapuscrit de Deligne, p. 16 : preuve de la proposition 6.9, (I) à (V). Sans annotation manuscrite. La preuve de (V) n'est pas achevée sur cette page : le tapuscrit se poursuit au-delà du lot.}

\page{20}
\note{Verso resté blanc, portant seulement une courte note de la main annotatrice à l'encre bleue (un corollaire pour $S$ plat sur $\mathbb{Z}$). Écriture d'un tiers : non transcrite.}

\end{document}
